\subsection{Properties of Equality} 
There are also some properties that apply when real numbers are equal.
\begin{displaybox}[Symmetric and Transitive Properties of Equality]
\begin{itemize}
\item For any real numbers $a$ and $b$, $a=b$ if and only if $b=a$.
\item For any real numbers $a$, $b,$ and $c$, if $a=b$ and $b=c$ then $a=c$.
\end{itemize}
\end{displaybox}
\begin{displaybox}[Additive Property of Equality]
For any real numbers $a$, $b$, and $c$, $a=b$ if and only if $a+c=b+c$.

We can think of this as ``adding'' or ``subtracting'' the same number on both sides. (Remember that subtracting is just adding the \makeblank[1.5in].)
\end{displaybox}
\begin{displaybox}[Multiplicative Property of Equality]
For any real numbers $a$, $b$, and $c$, with $c\neq 0$, $a=b$ if and only if $ac=bc$.

We can think of this as ``multiplying'' or ``dividing'' by the same number on both sides. (Remember that dividing is just multiplying by the \makeblank[1.5in].)
\end{displaybox}
\begin{example}
Assume $2x+2=4x-6$. Get equivalent equations by following each instruction in turn.
\begin{lpic}{}{9}
\regtextleft{0}{-2}{Add $4$ to both sides};
\regtextleft{0}{-4}{Multiply both sides by $3$};
\regtextleft{0}{-6}{Subtract $4x$ from both sides};
\regtextleft{0}{-8}{Divide both sides by $2$};
\end{lpic}
To check that these equations are all equivalent, we'll see that they're all true when $x=4$.
\begin{lpic}{}{5}
\end{lpic}
\end{example}